\documentclass[acmsmall,nonacm]{acmart}
\usepackage{minted}
\newcommand{\vtiny}{\fontsize{5.3}{5.8}\selectfont}
\renewcommand{\theFancyVerbLine}{%
  \fontsize{4.0}{4.0}\selectfont\arabic{FancyVerbLine}}
\newcommand{\blk}[2]{%
  \inputminted[fontsize=\vtiny,linenos,stepnumber=5,
    firstnumber=#2,xleftmargin=8pt,numbersep=5pt,
    frame=leftline,framerule=0.3pt,framesep=3pt,
    rulecolor=\color{black!25},breaklines]{python}{#1}}
\usepackage{wrapfig,needspace}
\makeatletter
\newcommand{\wfclear}{\par
  \ifnum\c@WF@wrappedlines>0
    \null\vskip\c@WF@wrappedlines\baselineskip
    \c@WF@wrappedlines\z@
  \fi}
\makeatother
\newenvironment{codepage}[2]
  {\needspace{#1}\def\capt{#2}%
   \begin{wrapfigure}{R}{0.55\textwidth}}
  {\captionof{listing}{\capt}\end{wrapfigure}
   {\sloppy\prose\par}\wfclear\medskip}
\title[ezr.py]{Can AI Be Small? A Complete, Explainable,
Multi-Objective Toolkit in 444 Lines}
\author{Tim Menzies}
\affiliation{\institution{North Carolina State University}
  \city{Raleigh}\country{USA}}
\email{timm@ieee.org}
\begin{document}
\begin{abstract}
Much recent commentary asserts that effective AI requires
large models, large data, and large infrastructure. We offer
a counter-example, at least for tabular software-engineering
optimization tasks where independent attributes are cheap to
read but dependent goals are expensive to measure. This
paper presents, in its entirety, ezr.py: 444 lines of
dependency-free Python implementing incremental column
summaries, multi-objective distance, active learning,
Naive Bayes classification, classification and regression
trees, a confusion-matrix scorer, and the non-parametric
statistics needed to check all of the above. On the 120+
tasks of the MOOT repository, this code's active learner
finds near-optimal rows after labelling around fifty
examples, matching heavyweight optimizers that run orders
of magnitude slower. Unusually, the paper's format matches
its thesis: every line of the toolkit appears here, one
conceptual unit per page, with commentary alongside. We
argue that when a complete system fits in a paper's
figures, reading code becomes a research method: the
simplifications reported here were found not by adding
machinery, but by deleting it.
\end{abstract}
\maketitle

\section{Introduction}
It is now routinely argued that humans no longer need to
read or write code; that the era of hand-crafted software is
ending; that AI is, by nature, big. Whatever the merits of
those claims for perception or generation tasks, they are
testable in at least one domain: tabular software-engineering
optimization, where each row's independent attributes cost
nothing to read, but each row's goals (runtimes, defect
counts, build outcomes) cost a measurement. In that domain,
this paper defends a simple thesis: \emph{the AI that
matters fits on a page}---or, more precisely, on eight of
them.

The vehicle for that argument is ezr.py, a rewrite of our
earlier EZR toolkit that emerged from a months-long exercise
in deletion. Version by version (three complete rewrites,
each verified output-identical on 120+ datasets before being
trusted), classes collapsed into tuples and dicts, separate
add/subtract/summarize routines collapsed into one update
primitive, and separate learners collapsed into short
questions asked of shared summaries. The result is 444
lines, zero dependencies, in which:
\begin{itemize}
\item a \emph{Num} is a tuple $(n, \mu, m_2)$ and a
  \emph{Sym} is a dict of counts (Listing~2);
\item one primitive, add (whose inc=$-1$ undoes), maintains
  both, so rows can move between summaries in constant time;
\item active learning labels ${\sim}50$ of thousands of
  rows, then a small tree ranks the rest (Listings~3, 5);
\item Naive Bayes classification and its confusion-matrix
  scoring share the same columns (Listing~4);
\item effect-size and distribution statistics---Cliff's
  delta, Kolmogorov--Smirnov, Cohen's $d$---check every
  comparison the toolkit makes about itself (Listing~6).
\end{itemize}

Our contribution is not any single algorithm; all are
decades old. It is the demonstration, in full, that they
share so much structure that a complete, statistically
self-checking, explainable toolkit needs less code than
many papers' related-work sections. Prior work reported
this quantitatively: active learning on these summaries
matches state-of-the-art optimizers such as SMAC while
running orders of magnitude faster, and the trees it grows
explain themselves in under ten attributes even when
thousands are available. This paper adds the qualitative
case: we show every line, because at this size we can.

The paper's format follows its content. Section 2 and
onward present ezr.py one conceptual unit per page: the
listing (left) with commentary (right). Line numbers are
those of the released file (\url{github.com/timm/ezr};
\texttt{pip install ezr}); nothing is elided. Data comes from the MOOT
repository of multi-objective SE tasks. We close with the
threats and the invitation: before deploying a heavyweight
pipeline, run a small baseline---if the heavyweight cannot
beat these eight pages, it is technical debt.

\section{The Code}
Each of the following pages is one conceptual unit of
ezr.py. Read the listings in order and you have read the
entire system; read the right-hand columns and you have
read its design rationale. The units are: configuration;
structs; distance and acquisition; Bayes; trees; reports
and statistics; evaluation and demos; and the sixteen-line
harness that replaces a test framework.
\clearpage

\def\prose{Every setting lives once, in the docstring, as
\texttt{-Key=default} plus a comment. One regex turns that text
into the global \texttt{the}; \texttt{--help} prints the same
text. Config, help and documentation cannot drift, because they
are one artifact.

\texttt{atom} coerces strings once, at the border; everything
downstream trusts its types. Its ordering is load-bearing: we
tried a ``faster'' variant that returned only floats, and it was
both slower (532 vs.\ 497 cpu-seconds over 120+ datasets---most
cells are integers, so \texttt{int} succeeds first try either
way) and wrong (settings such as \texttt{Few=128} became floats
and crashed every slice that consumed them). A ``richer''
variant that stripped and tested for booleans on every cell cost
13\% of total runtime. The shipped compromise recognizes
booleans only on the rare path, after both numeric parses fail:
measured cost, under 3\%, in the noise.

The \texttt{defaults} snapshot supports a contract used
throughout: demos may mutate \texttt{the} freely, and the runner
restores pristine settings afterwards, so no demo can leak into
the next. The CSV reader is five lines: no dtype inference, no
chunked reads, no NA conversion---pandas offers all that, and
this paper needs none of it. The \texttt{\$MOOT} expansion means
per-machine data without per-machine edits.}
\begin{codepage}{370pt}{Configuration and input. The docstring is
the single source of truth for all settings.}
\blk{sec00.py}{1}
\end{codepage}

\def\prose{The whole substrate is two tiny types. A
\texttt{Num} is a tuple $(n,\mu,m_2)$---all that Welford's
algorithm needs to track a running mean and standard deviation,
in one pass, without catastrophic cancellation. A \texttt{Sym}
is a dict of counts. One primitive, \texttt{add}, updates
either; called with \texttt{inc=-1} it \emph{undoes}. That
symmetry is what makes data movement cheap everywhere else: the
tree's cut sweeper slides values between summaries with one add
each way, and the active learner demotes a row from its elite
pool in constant time. There is no retraining step anywhere in
this paper.

Two traps are worth recording. \texttt{adds} guards with
\texttt{if it is None}, not \texttt{it or Num()}: an empty
\texttt{Sym()} is falsy, and the shorter idiom silently replaces
a fresh symbol counter with a number summary. And code testing
``is this the first column?'' must write \texttt{is not None},
because column index zero is falsy too.

\texttt{Tbl} routes columns by header convention: upper-case
names are numeric; \texttt{+}/\texttt{-} suffixes mark goals to
maximize or minimize; \texttt{!} marks a class column;
\texttt{X} means ignore. Decisions made once, at header time,
propagate everywhere---no re-parsing of names later. One type
test (\texttt{type(col) is Sym}) plus functions-as-values buys
what class hierarchies buy, without the hierarchy; the pattern
is the same one the Linux kernel uses with structs of function
pointers, thirty years and no classes.}
\begin{codepage}{345pt}{Structs: Num, Sym, and the one update
primitive that both builds and unbuilds them.}
\blk{sec01.py}{54}
\end{codepage}

\def\prose{Numbers are normalized by a clamped z-score
pushed through a logistic. The clamp (at $\pm3$ standard
deviations) stops one outlier crushing everyone else onto a dot;
the 1.7 in the logistic makes it approximate the normal CDF, so
columns with wildly different units all speak the same $[0,1]$
language. \texttt{ydist} then scores a row by Minkowski distance
to ``heaven''---every goal at its best---and \texttt{xdist} is
the same norm over independent columns, with missing values
treated pessimistically (distance 1).

On top of distance, twenty lines give active learning. Keep two
pools: \emph{best} (sorted, elite) and \emph{rest}. Label next
whichever unlabelled row is nearest best's centroid and farthest
from rest's. After each label, if best has outgrown
$\sqrt{1+b+r}$, demote its worst row. The \texttt{1+} is not
decoration: it grants one row of hysteresis exactly when the
pools are tiny and a one-row ``best'' would make the centroid
pure noise. Deleting that single character shifted results on 76
of 127 benchmark datasets and cost a median point.

Budget honesty matters more than either. \texttt{acquire} takes
a cap, and the evaluation harness passes
\texttt{Stop$-$Check}---the labels reserved for checking the
model's suggestions are counted against the same budget as the
labels that built it. Instrumenting that honestly \emph{raised}
the observed median (84 to 87): verification labels, spent on
the model's top picks, out-earn acquisition labels.}
\begin{codepage}{365pt}{Distance and acquisition: normalize,
score against heaven, then label only the informative
rows.}
\blk{sec02.py}{103}
\medskip
\blk{sec03.py}{134}
\end{codepage}

\def\prose{Bayes in three functions: \texttt{like} is a
smoothed frequency (Syms) or a Gaussian pdf (Nums);
\texttt{likes} multiplies one row's evidence in log space, with
an $m$-estimate on each column and a $k$-smoothed prior on class
size; \texttt{liked} returns the most likely of several tables.
Because \texttt{Num} and \texttt{Sym} columns are already fitted
summaries, there is no separate training phase: \texttt{add}
fits, \texttt{like} reads. A classifier is just a dict of
per-class clones of the table, filled by \texttt{addRow}.

\texttt{confuse} turns (got, want) pairs into per-class
counts and then accuracy, recall (\emph{pd}), false alarm
(\emph{pf}) and precision. Two details: entries are created
lazily, on first sight of a class, which is safe only because
true negatives are derived at the end
($tn = n - tp - fn - fp$) rather than counted per pair; and
classes never predicted and never wanted simply do not appear,
which is correct---they would be all-$tn$ rows saying nothing.

Reporting \emph{both} pd and pf is a methodological point, not
a convenience. In our text-mining replication work, a
recommended preprocessing step looked fine on recall alone and
was exposed as counterproductive only when false alarm was
measured beside it. Cheap tools make such re-checks cheap; this
scorer is fourteen lines, and it grades every classifier
comparison in this paper (Listing~7) on the same splits.}
\begin{codepage}{255pt}{Naive Bayes and the confusion-matrix
scorer.}
\blk{sec04.py}{155}
\end{codepage}

\def\prose{Trees are recursive discretization.
\texttt{cutNum} sweeps sorted (x,y) pairs, sliding each value
from the right summary to the left with one \texttt{add} each
way and yielding a candidate cut only where the x value actually
changes. \texttt{cutSym} tries each symbol one-vs-rest with
fresh accumulators. (We considered the classic Fisher trick of
ordering symbols by mean y and cutting prefixes; for nominal
attributes one-vs-rest reads better as a rule, and at
active-learning sizes---tens of labelled rows---the asymptotic
saving is irrelevant.) Splits are scored by size-weighted
expected diversity: standard deviation for numbers, entropy for
symbols.

That accumulator choice is the whole classification story. The
tree peeks at the type of its first y value: numeric y (distance
to heaven) makes regression trees; symbolic y (a class column)
makes classification trees. One line changes; the substrate does
not. Missing values route by imputing the column's middle---mean
for numbers, mode for symbols; the symbolic half of that
symmetry was a late fix, and matters on datasets with missing
nominals.

One tuning note. \texttt{Leaf=4} is tuned for the active
learner, which grows trees from ${\sim}45$ labelled rows. Grown
passively from hundreds of rows, such trees overfit: on the
diabetes data, accuracy rose from 0.69 to 0.75 simply by scaling
the leaf minimum to $\sqrt{n}$. Small data and big data want
different regularization, and one knob expresses both.}
\begin{codepage}{395pt}{Trees: recursive discretization, one
accumulator swap away from classification.}
\blk{sec05.py}{189}
\end{codepage}

\def\prose{Reporting walks the tree, printing each leaf's
distance-to-heaven, size, and per-goal centroids, with
\texttt{+}/\texttt{-} marking the best and worst leaves. The
display \emph{is} the model: a manager reading it is reading the
actual decision boundary, not a post-hoc approximation of it,
which is the property that lets tiny trees stand in for SHAP- or
LIME-style explanation tools.

The stats section holds the three checks used everywhere in
this paper. Cliff's delta asks whether two samples' rank
imbalance exceeds 0.197 (a ``small'' effect); the
Kolmogorov--Smirnov statistic asks whether their distributions
diverge beyond the 95\% line ($1.36\sqrt{(n+m)/nm}$); Cohen's
$d$ asks whether the mean gap exceeds 0.35 pooled standard
deviations, with an optional absolute floor (\texttt{eps}) for
scales, like accuracy, where tiny differences are meaningless
regardless of variance. \texttt{same} declares two result
distributions indistinguishable only when all three agree.

This conservatism is deliberate. Pre-registering an effect-size
threshold and aggregating over twenty or more repeats rejects
trivially small differences without running per-cell hypothesis
tests across dozens of comparisons---the multiple-comparison
correction problem never arises because the per-cell tests are
never run. The whole apparatus is 26 lines, which is the reason
it actually gets used: every claim of the form ``X beats Y''
made by this toolkit's own demos passes through
\texttt{same} first.}
\begin{codepage}{330pt}{Reports and statistics: show the tree;
ask if two treatments really differ.}
\blk{sec06.py}{246}
\medskip
\blk{sec07.py}{264}
\end{codepage}

\def\prose{\texttt{wins} converts distance-to-heaven into
the complement of normalized regret: 100 means ``found the best
row in the table'', 0 means ``no better than the average row'',
negative means worse than that. Normalizing this way lets one
number compare treatments across a hundred datasets whose raw
goal scales differ by orders of magnitude.

\texttt{holdout} is the evaluation loop: shuffle; give the
active learner half the rows (capped at \texttt{Few}) minus the
\texttt{Check} labels reserved for verification; grow a tree
from what it labelled; then, on the untouched half, label only
the tree's top \texttt{Check} picks and keep the best. Twenty
repeats, report the mean. Over the 120+ MOOT optimization
tasks this loop scores a median win of 86--87, and the full
sweep runs in about a minute of wall-clock time on a laptop
(500 cpu-seconds at 10-way parallelism)---which is precisely
what makes the re-checking culture of the previous page
affordable.

The first demos are worked examples of the substrate: Welford
matching the textbook mean and standard deviation; the entropy
of \texttt{aabbbc}; header routing sending columns to x, y,
klass or nowhere; one legal, routable cut. Note that each demo
prints its evidence (\texttt{mu 5.0 sd 2.138}), not
``ok'': the demo teaches, the output is diffable, and a wrong
number stays visible even when an assert is too loose.}
\begin{codepage}{415pt}{Evaluation (wins, holdout) and the first
demos.}
\blk{sec08.py}{292}
\medskip
\blk{sec09.py}{308}
\end{codepage}

\def\prose{The larger demos. \texttt{test\_tree} grows and
shows a tree on the optimization data; \texttt{test\_holdout}
runs the twenty-repeat evaluation; \texttt{test\_klass} is a
classifier bake-off in which tree and bayes see the \emph{same}
thirty 50:50 splits (a paired design that removes split-to-split
variance from the comparison), their pooled confusions are
printed side by side, and the statistics of Listing~6 rule on
the difference. On the diabetes data the verdict reads:
bayes 75\% vs.\ tree 73\%, delta 2, \emph{different}. Naive
Bayes degrades gracefully with less training data; trees
fragment---exactly the regime distinction the \texttt{Leaf}
note on Listing~5 predicts. \texttt{\_klass} owns the
scaffolding; each classifier is a \texttt{fit} function
returning a predictor, so adding a third contender is five
lines.

Below them, the harness. \texttt{run} seeds the random number
generator before every demo, so any demo runs alone or in any
order, and restores pristine settings in a \texttt{finally}
(one trap: a \texttt{return} \emph{inside} finally would
silently swallow the crash count---the \texttt{return 0} must
sit outside). \texttt{cli} evaluates the argument list left to
right: \texttt{-Key val} sets, \texttt{--name} runs, anything
else warns and continues. The exit code counts crashes, so a
Makefile gates on the suite with no framework. Any function
named \texttt{test\_*} is a demo, a test, and a help entry all
at once: conventions are load-bearing.}
\begin{codepage}{560pt}{The larger demos and the harness:
a statistically-checked bake-off, then sixteen lines that
replace a test framework.}
\blk{sec10.py}{352}
\medskip
\blk{sec11.py}{426}
\end{codepage}



\end{document}
